% Part: first-order-logic
% Chapter: introduction
% Section: soundness-completeness

\documentclass[../../../include/open-logic-section]{subfiles}

\begin{document}

\olfileid[pt]{fol}{int}{scp}

\olsection{Correção e Completude}

Também introduziremos sistemas de !!{derivation} para a lógica de
primeira ordem.  Há muitos sistemas de !!{derivation} que os lógicos
desenvolveram, mas todos definem a mesma relação de !!{derivability}
entre !!{sentence}s. Dizemos que $\Gamma$ \emph{!!{derive}}~$!A$,
$\Gamma \Proves !A$, se há uma !!{derivation} de um certo tipo
precisamente definido. As !!{derivation}s são sempre disposições
finitas de símbolos---talvez uma lista de !!{sentence}s, ou alguma
estrutura mais complicada.  O propósito dos sistemas de
!!{derivation} é fornecer uma ferramenta para determinar se uma
!!{sentence} é consequência de algum conjunto~$\Gamma$.  Para servir a
esse propósito, deve ser verdade que $\Gamma \Entails !A$ se, e
somente se, $\Gamma \Proves !A$.

Se $\Gamma \Proves !A$ mas não $\Gamma \Entails !A$, nosso sistema de
!!{derivation} seria forte demais, provaria demais.  A propriedade de
que se $\Gamma \Proves !A$ então $\Gamma \Entails !A$ é chamada de
\emph{correção}, e é um requisito mínimo para qualquer bom sistema de
!!{derivation}. Por outro lado, se $\Gamma \Entails !A$ mas não
$\Gamma \Proves !A$, então nosso sistema de !!{derivation} é fraco
demais, não prova o suficiente.  A propriedade de que se $\Gamma
\Entails !A$ então $\Gamma \Proves !A$ é chamada de
\emph{completude}. A correção é usualmente relativamente fácil de
provar (por indução na estrutura das !!{derivation}s, que são
definidas indutivamente). A completude é mais difícil de provar.

Correção e completude têm várias consequências importantes. Se um
conjunto de !!{sentence}s~$\Gamma$ !!{derive} uma contradição (como
$!A \land \lnot !A$), ele é chamado de \emph{inconsistente}.
$\Gamma$s inconsistentes não podem ter modelos; são insatisfatíveis. Da
completude segue a contrapositiva: qualquer $\Gamma$ que não seja
inconsistente---ou, como diremos, \emph{consistente}---tem um modelo.
De fato, isso é equivalente à completude, e é a forma de completude que
de fato provaremos.  É um resultado profundo e talvez surpreendente:
apenas o fato de que você não pode provar $!A \land \lnot !A$ a partir
de $\Gamma$ garante que há uma !!{structure} que é como $\Gamma$ a
descreve.  Assim, a completude dá uma resposta à pergunta: quais
conjuntos de !!{sentence}s têm modelos? Resposta: todos e somente os
conjuntos consistentes têm.

Os teoremas de correção e completude têm duas consequências
importantes: a compacidade e o teorema de L\"owenheim--Skolem. Estes
são resultados importantes na teoria de modelos e podem ser usados
para estabelecer muitos resultados interessantes. Já mencionamos dois:
a lógica de primeira ordem não pode expressar que o !!{domain} de uma
!!{structure} é finito ou que é !!{nonenumerable}.

Historicamente, tudo isso---como definir a sintaxe e a semântica da
lógica de primeira ordem, como definir bons sistemas de
!!{derivation}, como provar que são corretos e completos, ficar claro
sobre o que pode e não pode ser expresso em linguagens de primeira
ordem---levou muito tempo para ser descoberto e acertado.  Agora
sabemos como fazer isso, mas percorrer todos os detalhes ainda pode ser
confuso e tedioso.  Mas também é importante, porque os métodos
desenvolvidos aqui para a linguagem formal da lógica de primeira ordem
são aplicados em todo lugar em lógica, ciência da computação e
linguística. Assim, trabalhar os detalhes compensa a longo prazo.

\end{document}
